% CLASS 4, part two -- the notes' section 1.5, sec.lobsterEmpirics.
% The chapter's assignment, briefed and set. Planned 18 min over 5 frames.
% Under \includeonly this file builds alone, so it carries its own Outline.
%
% THE FRAME CONTRACT. These slides are projection material, not reference material: the
% students read the notes, and they only ever see a frame while it is being talked over.
% So every frame carries ONE OBJECT -- a display, a statement of at most two lines, a TikZ
% picture, a table or an algorithmic block -- and under it at most three CAPTIONS. A caption
% is a gloss (`term: noun phrase'), a fragment of subject and predicate, or a condition
% written as mathematics. It is never a sentence, never a subordinate clause, and never a
% repetition of the frame title: what the title says, the body does not say again, and what
% follows from the object is spoken. At most 12 words of prose each, 25 to a frame, hard
% cap 40. A statement states its head claim and stops; its consequence is spoken.
% What is said over the object goes in \note{}, which \documentclass[script]{unito26slides}
% prints beside each frame and the deck itself never shows. Check with
%   python3 .claude/skills/writing-in-tex/scripts/frame_density.py <file>
% Open the class on the chapter roadmap, as classes 1 and 2 do. Do not close on a summary
% frame: a slide that recites what the class just did is the spoken half written down.
%
% LABELS. eq.queueImbalance belongs to class 1 and remark.passiveDirection to class 2; both
% are RESTATED HERE WITHOUT LABELS. Nobody remembers a formula from three weeks ago, and the
% assignment turns on both of them.

\subsection{Assignment --- evidence from recorded data}

\begin{frame}
  \frametitle{Outline}
% 2 min. One \nocite per line -- a multi-key \nocite spanning lines breaks bibtex.
\nocite{LPV22sho}
\nocite{HP11lob}
% The apparatus is developed on flow we generate, and a result obtained on generated flow is
% a result about the generator.
\end{frame}

\begin{frame}
  \frametitle{The intensities do not read the book}
% 4 min. Redisplay eq.queueImbalance here, unlabelled. A resting limit order is short an
% option: it is filled when someone wishes to trade against it, which is disproportionately
% when trading against it is right, and the defence is to withdraw first. So a queue about to
% be picked off is one whose own side pulls away from it, and depletion begets depletion.
% That is what makes \queueImbalance[1] informative in a traded book -- a thin bid is thin
% BECAUSE those posting on that side inferred something. A kernel built for replenishment
% does the reverse, and the two cannot both be strong in one specification.
\end{frame}

\begin{frame}
  \frametitle{Three more absences}
% 3 min. The marks are independent of the flow, and a generator whose withdrawals draw
% against the size actually resting emits a THINNED process rather than the Hawkes process it
% started from -- the hypothesis fails first on the withdrawal types. One excitation
% timescale; no intraday non-stationarity; and no anchor for the price level, so the mean
% price change over a session is a property of the seed.
\end{frame}

\begin{frame}
  \frametitle{Read the comparison asymmetrically}
% 4 min. remark.asymmetricReading. Whatever forecasting power \queueImbalance[1] retains in
% the generated setting is MECHANICAL -- a depleted queue is one market order from clearing
% -- and none of it is informational.
% A result favourable to \queueImbalance[1] is the stronger finding.
% Then what recorded data settles: whether \signedExcitation carries either sign in a traded
% market, and over what window and horizon the sign of the forecast can be read at all.
\end{frame}
